\chapter{\ifenglish Background Knowledge and Theory\else ทฤษฎีที่เกี่ยวข้อง\fi}

เนื้อหาในบทนี้อธิบายแนวคิดและเทคโนโลยีที่เกี่ยวข้องกับโครงงาน
ได้แก่ ระบบ SIEM แหล่งข้อมูล Logs มาตรฐานการจัดรูปแบบข้อมูลและ Detection Rule กรอบความรู้ MITRE ATT\&CK
ชุดข้อมูลที่ใช้ทดลอง โมเดลภาษาขนาดใหญ่ เทคนิค RAG รูปแบบการนำเสนอข้อมูล
ความสัมพันธ์เชิงกราฟของ Process และตัวชี้วัดที่ใช้ประเมินผล
เพื่อเป็นพื้นฐานสำหรับการทำความเข้าใจการออกแบบในบทถัดไป

\section{Security Information and Event Management (SIEM)}

SIEM คือระบบที่รวบรวม Logs และเหตุการณ์ด้านความปลอดภัย (Security Events)
จากแหล่งต่าง ๆ ภายในองค์กร เช่น เครื่องลูกข่าย เซิร์ฟเวอร์ ไฟร์วอลล์ และแอปพลิเคชัน
มาจัดเก็บไว้ในที่เดียว เพื่อให้สามารถค้นหา เชื่อมโยง (Correlate) และสร้างการแจ้งเตือน (Alert)
เมื่อพบพฤติกรรมที่ตรงกับกฎการตรวจจับ (Detection Rule)
SIEM จึงเป็นเครื่องมือหลักของศูนย์ปฏิบัติการด้านความปลอดภัย (Security Operations Center: SOC)

ในโครงงานนี้ใช้ \textbf{Splunk Enterprise} เป็นระบบ SIEM
โดย Logs ถูกจัดเก็บเป็นดัชนี (Index) และสามารถเรียกค้นได้ด้วยภาษา Search Processing Language (SPL)
หรือผ่าน REST API ซึ่ง Pipeline ของโครงงานใช้สำหรับดึง Logs ที่เกี่ยวข้องกับ Alert

\subsection{Detection Rule และ ES Content Update (ESCU)}
Alert ใน Splunk เกิดจาก Detection Rule ซึ่งเป็นคำค้น SPL ที่ทำงานตามเงื่อนไข
Splunk เผยแพร่ชุด Detection Rule ทางการในชื่อ ES Content Update (ESCU)~\cite{splunk-escu}
ซึ่งติดตั้งเป็น App เดียวและระบุเทคนิค ATT\&CK ที่แต่ละ Rule ตรวจจับ
Rule แบ่งเป็น 3 ประเภท ได้แก่ \textbf{TTP} (พฤติกรรมการโจมตีชัดเจน สร้าง Alert โดยตรง)
\textbf{Anomaly} (พฤติกรรมผิดปกติ ส่งคะแนนเข้าระบบความเสี่ยง) และ
\textbf{Hunting} (ใช้ค้นหาด้วยมือ ไม่สร้าง Alert)
Alert ที่เกิดจากการโจมตีจริงเรียกว่า True Positive (TP) ส่วน Alert ที่เกิดจากกิจกรรมปกติเรียกว่า False Positive (FP)
ในโครงงานนี้ใช้ Rule ประเภท TTP และ Anomaly ที่ใช้ข้อมูล Windows Sysmon รวม 829 Rule
สร้าง Alert จาก Logs ที่ Replay เข้า Splunk เพื่อจำลองสถานการณ์ที่ SOC Analyst ได้รับ Alert จริง

\subsection{SOC Analyst}
SOC Analyst คือเจ้าหน้าที่ด้านความปลอดภัยไซเบอร์ที่ทำหน้าที่เฝ้าระวัง ตรวจสอบ คัดกรอง
และวิเคราะห์เหตุการณ์ด้านความปลอดภัย เพื่อระบุภัยคุกคามและสนับสนุนการตอบสนองต่อเหตุการณ์ (Incident Response)
ซึ่งเป็นผู้มีส่วนได้ส่วนเสีย (Stakeholder) หลักของโครงงานนี้

\section{ชุดข้อมูลและแนวคิด Scenario Dataset}

\subsection{Splunk Attack Data และ Atomic Red Team}
Splunk Attack Data~\cite{splunk-attack-data} เป็นคลังชุดข้อมูลของ Splunk Threat Research Team
ที่บันทึก Logs จากการจำลองการโจมตีใน Lab ที่สร้างด้วยเครื่องมือ attack\_range
แต่ละชุดมีไฟล์ YAML ระบุเทคนิค ATT\&CK ที่จำลอง แต่ไม่ได้ระบุว่า Event ใดเป็นส่วนของการโจมตี
ชุดข้อมูลบางส่วนสร้างด้วย Atomic Red Team (ART)~\cite{atomic-red-team}
ซึ่งเป็นคลังชุดทดสอบการโจมตีขนาดเล็กที่แยกตามเทคนิคและเขียนคำสั่งไว้ชัดเจน
และคำสั่ง \texttt{Invoke-AtomicTest} ที่ใช้รัน หรือไฟล์จากโฟลเดอร์ \texttt{AtomicRedTeam}
จะปรากฏใน Logs เป็นจุดเริ่มต้นของการโจมตี
จึงสามารถระบุ Event ที่เป็นการโจมตีได้อัตโนมัติจาก Process Tree
โครงงานจึงใช้ชุดข้อมูล ART ของ Splunk เป็นชุดข้อมูลสำหรับพัฒนา (Development Set)

\subsection{OTRF Security-Datasets}
OTRF Security-Datasets~\cite{otrf-security-datasets} (เดิมชื่อ Mordor) เป็นชุดข้อมูลสาธารณะของ Open Threat Research Forge
ที่บันทึก Logs จากการจำลองการโจมตีใน Lab ของตนเอง โดยใช้เครื่องมืออย่าง Empire และ Covenant
แบ่งเป็นชุดแบบ Atomic (จำลองพฤติกรรมเดียว) และ Compound (จำลองทั้ง Campaign ที่มีหลายเทคนิค)
และมี Metadata ส่วน Adversary View ที่บันทึกคำสั่งที่ผู้โจมตีใช้จริง
เนื่องจากมาจาก Lab และเครื่องมือที่ต่างจาก Splunk จึงใช้เป็นชุดข้อมูลสำหรับทดสอบ (Test Set)
เพื่อตรวจสอบว่าค่าที่ปรับบนข้อมูลของ Splunk ใช้กับข้อมูลจากแหล่งอื่นได้หรือไม่
โครงงานใช้เฉพาะชุดข้อมูลประเภท Atomic เนื่องจากแต่ละชุดมีเทคนิคที่เป็นเฉลยเพียงเทคนิคเดียว
ส่วนชุดประเภท Compound มีหลายเทคนิคปนกันในสถานการณ์เดียว จึงไม่สามารถใช้เป็น Ground Truth ระดับ Scenario ได้

\subsection{นิยามที่ใช้ในโครงงาน}
\begin{itemize}
  \item \textbf{Attack Scenario} -- การจำลองการโจมตีหนึ่งสถานการณ์ในสภาพแวดล้อมการทดลอง
  \item \textbf{Scenario Dataset} -- ชุดของ Events ที่เกิดขึ้นในบริบทของ Attack Scenario หนึ่ง
  \item \textbf{Atomic Dataset} -- Scenario Dataset ที่จำลองการโจมตีเพียง 1 เทคนิค
      (1 Scenario = 1 เทคนิค) ตรงข้ามกับแบบ Compound ที่จำลองหลายเทคนิคต่อเนื่องกัน
  \item \textbf{Event / Log Event} -- บันทึกหนึ่ง Record ของกิจกรรม เช่น Process Creation หรือ Network Connection
  \item \textbf{Ground Truth} -- เทคนิค ATT\&CK ที่ใช้สร้าง Scenario นั้น
        เป็น Ground Truth ระดับ Scenario/Dataset ไม่ใช่ Label ของทุก Event
        (1 Scenario Dataset $\rightarrow$ 1 Ground-truth Technique $\rightarrow$ หลาย Events)
  \item \textbf{Ground Truth ระดับ Event} -- รายการ Events ที่เป็นส่วนของการโจมตีจริง
        มีเฉพาะใน Development Set ใช้ปรับพารามิเตอร์และวัดสัดส่วนหลักฐานที่ถูกคัดเลือก
  \item \textbf{Case (กรณีการสืบสวน)} -- หน่วยของการทดลอง ประกอบด้วย 1 Alert และ Logs ของไฟล์นั้น
  \item \textbf{Anchor} -- Event ที่ Alert ชี้ เป็นจุดเริ่มต้นของ Event Selection
        หากไฟล์ไม่มี Alert จะใช้ Event ของคำสั่งโจมตีเป็น Anchor สำรอง
  \item \textbf{Development Set / Test Set} -- ชุดข้อมูลสำหรับปรับพารามิเตอร์ (ใช้ได้หลายครั้ง)
        และชุดข้อมูลสำหรับวัดผล (ใช้ครั้งเดียวหลังล็อกค่า)
  \item \textbf{Freeze} -- การล็อกค่าพารามิเตอร์ Prompt เวอร์ชันของชุดข้อมูล และ Checksum ก่อนทดสอบ
  \item \textbf{Leak Check} -- การตรวจสอบว่า Prompt ไม่มีเฉลยหลุดเข้าไป เช่น รหัสเทคนิค หรือชื่อ Detection Rule
  \item \textbf{Noise} -- เหตุการณ์ใน Logs ที่ไม่เกี่ยวข้องกับการวิเคราะห์ รวมถึง Field และ Event ที่ไม่จำเป็น
  \item \textbf{Investigation Context} -- ชุดเหตุการณ์และความรู้ที่เกี่ยวข้องกับเหตุการณ์ที่น่าสงสัย
        ซึ่งถูกจัดเตรียมเพื่อให้ LLM ใช้เหตุผลกับเหตุการณ์นั้น
\end{itemize}

\section{ประเภท Logs ในชุดข้อมูล}
Logs ในชุดข้อมูลทั้งสองแหล่งถูกบันทึกจากเครื่องในสภาพแวดล้อมจำลองของผู้จัดทำชุดข้อมูล
โดยใช้เครื่องมือบันทึกเหตุการณ์ของระบบปฏิบัติการ ดังนี้

\subsection{Windows Sysmon}
System Monitor (Sysmon)~\cite{sysmon} เป็นบริการของ Microsoft Sysinternals
ที่บันทึกกิจกรรมของระบบปฏิบัติการ Windows อย่างละเอียดลงใน Windows Event Log
โดยแต่ละประเภทกิจกรรมมีหมายเลข Event ID เฉพาะ ตัวอย่าง Event ID ที่สำคัญต่อโครงงานแสดงในตารางที่~\ref{tab:sysmon-eid}
Events ที่เกี่ยวข้องกับ Process จะมีตัวระบุ \texttt{ProcessGuid} ที่ไม่ซ้ำกัน
และ Event ID 1 จะบันทึก \texttt{ParentProcessGuid} ไว้ด้วย
จึงสามารถนำมาสร้างความสัมพันธ์แบบ Parent--Child ระหว่าง Process ได้

\begin{table}[ht]
  \centering
  \caption{ตัวอย่าง Sysmon Event ID ที่ใช้ในโครงงาน}
  \label{tab:sysmon-eid}
  \begin{tabular}{|c|l|l|}
    \hline
    \textbf{Event ID} & \textbf{ชื่อเหตุการณ์} & \textbf{ข้อมูลสำคัญ} \\ \hline
    1  & Process Creation        & Image, CommandLine, ParentImage, User \\ \hline
    3  & Network Connection      & SourceIp, DestinationIp, DestinationPort \\ \hline
    10 & Process Access          & SourceImage, TargetImage, GrantedAccess \\ \hline
    11 & File Create             & Image, TargetFilename \\ \hline
    23 & File Delete (archived)  & Image, TargetFilename \\ \hline
  \end{tabular}
\end{table}

\subsection{Linux Auditd}
Linux Audit Daemon (auditd)~\cite{auditd} เป็นระบบบันทึกเหตุการณ์ของ Linux Kernel
เช่น การเรียกใช้ System Call การรันคำสั่ง และการเข้าถึงไฟล์
โดยหนึ่งเหตุการณ์อาจประกอบด้วยหลาย Record (เช่น \texttt{SYSCALL}, \texttt{EXECVE}, \texttt{PATH})
ที่เชื่อมกันด้วยหมายเลข Event เดียวกัน
ในโครงงานนี้ Auditd อยู่ในขอบเขตการออกแบบ ส่วนการทดลองใช้ Windows Sysmon

\section{Common Information Model (CIM)}
เนื่องจาก Logs จากแต่ละแหล่งมีรูปแบบและชื่อ Field แตกต่างกัน
(เช่น Sysmon เป็น XML ส่วน Auditd เป็นข้อความแบบ Key=Value)
Splunk จึงกำหนด Common Information Model (CIM)~\cite{splunk-cim}
ซึ่งเป็นมาตรฐานชื่อ Field และโครงสร้างข้อมูลร่วม เช่น \texttt{process}, \texttt{parent\_process}, \texttt{user}, \texttt{dest}
% การ Normalize Logs ให้อยู่ในรูปแบบร่วมทำให้ขั้นตอนถัดไปของ Pipeline ทำงานได้กับ Logs หลายแหล่งโดยไม่ต้องเขียนตรรกะแยก
ภายใน Splunk การแปลง Logs ดิบให้เข้ากับ CIM ทำโดย Technology Add-on (TA) เช่น Splunk Add-on for Sysmon
และ Detection Rule ส่วนใหญ่ค้นผ่าน Data Model ของ CIM เช่น Rule ของ ESCU ที่ใช้ Sysmon 535 จาก 928 Rule
ค้นผ่าน Data Model \texttt{Endpoint.Processes}
ขั้นตอน Normalize ของโครงงานจึงแปลง Logs ทุกรูปแบบให้เป็น JSON ที่ใช้ชื่อ Field ตาม CIM
เพื่อให้ขั้นตอนถัดไปของ Pipeline ทำงานกับ Logs หลายแหล่งได้ด้วยตรรกะเดียวกัน

\section{MITRE ATT\&CK}
MITRE ATT\&CK~\cite{mitre-attack,strom2018attack} เป็นฐานความรู้ระดับสากลที่รวบรวมพฤติกรรมของผู้โจมตีไซเบอร์จากเหตุการณ์จริง
โดยจำแนกเป็น \textbf{Tactic} (เป้าหมายของการโจมตี เช่น Discovery, Credential Access)
\textbf{Technique} (วิธีการที่ใช้เพื่อบรรลุเป้าหมาย) และ \textbf{Sub-technique} (รายละเอียดเฉพาะของวิธีการ)
แต่ละเทคนิคมีรหัสประจำ เช่น T1087 Account Discovery และ T1087.001 Account Discovery: Local Account~\cite{mitre-t1087-001}
ซึ่งผู้โจมตีพยายามดึงรายชื่อบัญชีผู้ใช้ภายในเครื่อง เช่น ใช้คำสั่ง \texttt{net user} หรือ \texttt{net localgroup}

ATT\&CK Enterprise เวอร์ชัน 19.2 ที่ใช้ในโครงงานประกอบด้วย 15 Tactics, 222 Techniques
และ 475 Sub-techniques (ไม่นับรายการที่ถูกยกเลิก) และมีการปรับปรุงเวอร์ชันประมาณปีละ 2 ครั้ง
ในโครงงานนี้ ATT\&CK ถูกใช้ใน 2 บทบาท คือ (1) เป็นป้ายกำกับ (Label) ของ Ground Truth สำหรับประเมินผล
และ (2) เป็นฐานความรู้ (Knowledge Base) ให้ LLM อ้างอิงผ่าน RAG

\section{Large Language Model (LLM)}
LLM คือโมเดลภาษาที่สร้างจากสถาปัตยกรรม Transformer~\cite{vaswani2017attention}
และฝึกด้วยข้อความจำนวนมหาศาล ทำให้สามารถเข้าใจและสร้างข้อความ รวมถึงใช้เหตุผลกับข้อมูลที่ได้รับได้
LLM ประมวลผลข้อความในหน่วย \textbf{Token} และรับข้อมูลได้จำกัดตามขนาด \textbf{Context Window}
ต้นทุนและเวลาในการประมวลผลจึงแปรผันตามจำนวน Token ที่ส่งเข้าไป
นอกจากนี้ งานวิจัยพบว่าเมื่อ Context ยาวขึ้น LLM มีแนวโน้มใช้ข้อมูลที่อยู่ช่วงกลางของ Context ได้ไม่ดี~\cite{liu2024lost}
การคัดเลือกและจัดระเบียบข้อมูลก่อนส่งให้ LLM จึงช่วยลด Noise และทำให้กลไก Attention ของโมเดลโฟกัสข้อมูลที่สำคัญได้ดีขึ้น

\subsection{Local LLM: Qwen3.5 และ Ollama}
โครงงานเลือกใช้ LLM ขนาดเล็กที่ทำงานภายในเครื่อง ได้แก่ Qwen3.5 ขนาด 9B พารามิเตอร์
ซึ่งเป็นโมเดลในตระกูล Qwen~\cite{qwen3} ที่เผยแพร่แบบ Open-weight
และรันผ่าน Ollama~\cite{ollama} ที่ให้บริการโมเดลผ่าน API ภายในเครื่อง
ข้อดีของ Local LLM คือ Logs ขององค์กรไม่ต้องถูกส่งออกไปภายนอก และไม่มีค่าบริการ API ตามจำนวน Token
แต่โดยทั่วไปมีความสามารถจำกัดกว่าโมเดลเชิงพาณิชย์ขนาดใหญ่
เป้าหมายของโครงงานจึงเป็นการแสดงว่า Context ที่ดีช่วยให้โมเดลขนาดเล็กตอบได้ถูกต้อง

\section{Retrieval-Augmented Generation (RAG)}
RAG~\cite{lewis2020rag} คือเทคนิคที่เสริมความรู้ให้ LLM ขณะตอบคำถาม
โดยแบ่งเอกสารความรู้เป็นส่วนย่อย (Chunk) แล้วแปลงแต่ละ Chunk เป็นเวกเตอร์ด้วย Embedding Model
เก็บไว้ในฐานข้อมูลเวกเตอร์ (Vector Database)
เมื่อมีคำถาม (Query) ระบบจะแปลง Query เป็นเวกเตอร์และค้นหา Chunk ที่มีความคล้ายสูงสุด $k$ อันดับ (Top-$k$)
ด้วยการวัดความคล้ายแบบ Cosine Similarity ดังสมการ~\eqref{eq:cosine}
แล้วนำ Chunk เหล่านั้นใส่ใน Prompt เป็น Context ให้ LLM

\begin{equation}\label{eq:cosine}
\mathrm{sim}(\mathbf{q},\mathbf{d}) = \frac{\mathbf{q}\cdot\mathbf{d}}{\lVert\mathbf{q}\rVert\,\lVert\mathbf{d}\rVert}
\end{equation}

\noindent
โดยที่ $\mathbf{q}$ คือเวกเตอร์ของ Query และ $\mathbf{d}$ คือเวกเตอร์ของ Chunk ในฐานความรู้

\section{รูปแบบการนำเสนอข้อมูล (Data Representation)}
ข้อมูลชุดเดียวกันสามารถเขียนได้หลายรูปแบบ ซึ่งส่งผลต่อจำนวน Token และความเข้าใจของ LLM
รูปแบบที่โครงงานพิจารณา ได้แก่
\textbf{XML} (รูปแบบต้นฉบับของ Sysmon มีแท็กซ้ำซ้อนจำนวนมาก),
\textbf{JSON} (โครงสร้างชัดเจนและเป็น Output ของขั้นตอน Normalize),
\textbf{TOON} (Token-Oriented Object Notation~\cite{toon-format} ที่ตัดเครื่องหมาย \{\}, "" ออกเพื่อลด Token),
\textbf{Key=Value} (รูปแบบที่ Splunk ใช้แสดงผล) และ
\textbf{Timeline} (สรุปเหตุการณ์เป็นบรรทัดเดียวต่อ Event ใช้ Token น้อยที่สุดแต่ข้อมูลบางส่วนหายไป)
นอกจากนี้ยังมีรูปแบบ \textbf{Summary} ซึ่งให้ LLM แปลง Logs เป็นคำอธิบายพฤติกรรมภาษาธรรมชาติก่อน โดยไม่ระบุชื่อหรือรหัสเทคนิค เพื่อใช้เป็น Input หรือ Query สำหรับค้นหาในฐานความรู้

% \section{ความสัมพันธ์เชิงกราฟของ Process}
% ความสัมพันธ์ระหว่าง Process สามารถแทนด้วยกราฟมีทิศทาง $G=(V,E)$
% โดยแต่ละโหนด $v \in V$ คือ Process (ระบุด้วย ProcessGuid)
% และเส้นเชื่อม $e \in E$ แทนความสัมพันธ์ เช่น Parent สร้าง Child (Sysmon EID 1)
% หรือ Process หนึ่งเข้าถึงอีก Process (EID 10)
% การท่องกราฟ (Graph Traversal) แบบ Breadth-First Search ที่จำกัดจำนวนขั้น (Hop)
% ทำให้สามารถรวบรวม Events ที่เชื่อมโยงกับ Alert ได้โดยไม่ลามไปทั้งระบบ

\section{ความสัมพันธ์เชิงกราฟของ Process}
ความสัมพันธ์ระหว่าง Process สามารถแทนด้วยกราฟมีทิศทาง $G=(V,E)$
โดยแต่ละโหนด $v \in V$ คือ Process (ระบุด้วย ProcessGuid)
และเส้นเชื่อม $e \in E$ แทนความสัมพันธ์เชิงสาเหตุ ได้แก่
Parent สร้าง Child (Sysmon EID 1) และ Process หนึ่งสร้างไฟล์ที่ถูกรันเป็นอีก Process (EID 11 ตามด้วย EID 1)
ส่วน Process Access (EID 10) ใช้เป็นหลักฐานประกอบของโหนดที่เกี่ยวข้อง แต่ไม่ใช้เป็นเส้นทางในการท่องกราฟ
การท่องกราฟโดยหาเส้นทางสั้นที่สุด (Shortest Path) จากจุดที่ Alert ชี้และจำกัดระยะทางสูงสุด
ทำให้สามารถรวบรวม Events ที่เชื่อมโยงกับ Alert ได้โดยไม่ลามไปทั้งระบบ

\section{ตัวชี้วัดที่ใช้ประเมินผล}
\label{sec:metrics}

หน่วยของการประเมินผลคือกรณีการสืบสวน (Case) ไม่ใช่ราย Event เนื่องจาก Ground Truth อยู่ในระดับ Scenario Dataset
กำหนดให้ $N$ คือจำนวนกรณีทั้งหมด $y_i$ คือเทคนิคที่เป็นเฉลยของกรณีที่ $i$
และ $\hat{y}_i$ คือเทคนิคที่ระบบระบุ โดยใช้ตัวชี้วัดดังนี้

\begin{itemize}
  \item \textbf{Exact Match Accuracy} -- สัดส่วนกรณีที่ระบบระบุรหัสเทคนิคตรงกับเฉลยทุกหลัก
        รวมถึง Sub-technique
        \begin{equation}
        \mathrm{Acc}_{\mathrm{exact}} = \frac{1}{N}\sum_{i=1}^{N} \mathbf{1}\left[\hat{y}_i = y_i\right]
        \end{equation}

  \item \textbf{Parent Match Accuracy} -- สัดส่วนกรณีที่ระบบระบุถูกในระดับ Technique หลัก
        (เช่น ตอบ T1003.002 เมื่อเฉลยคือ T1003.001 นับว่าถูกในระดับ T1003)
        \begin{equation}
        \mathrm{Acc}_{\mathrm{parent}} = \frac{1}{N}\sum_{i=1}^{N} \mathbf{1}\left[\mathrm{parent}(\hat{y}_i) = \mathrm{parent}(y_i)\right]
        \end{equation}

  \item \textbf{RAG hit@$k$} -- สัดส่วนกรณีที่เทคนิคเฉลยอยู่ในรายการ Top-$k$ ที่ RAG ค้นคืนได้
        ใช้แยกปัญหาของขั้นตอนค้นคืน (Retrieval) ออกจากปัญหาของ LLM
        \begin{equation}
        \mathrm{hit@}k = \frac{1}{N}\sum_{i=1}^{N} \mathbf{1}\left[y_i \in R_k(q_i)\right]
        \end{equation}
        โดยที่ $R_k(q_i)$ คือรายการเทคนิค $k$ อันดับแรกที่ค้นคืนได้จาก Query $q_i$

  \item \textbf{Evidence Recall} -- สัดส่วนของ Events ที่เป็นการโจมตี (ตาม Ground Truth ระดับ Event)
        ที่ยังคงอยู่หลังขั้นตอน Event Selection วัดเฉพาะใน Development Set
        \begin{equation}
        \mathrm{Recall}_{\mathrm{evidence}} = \frac{|S \cap A|}{|A|}
        \end{equation}
        โดยที่ $S$ คือเซตของ Events ที่ถูกคัดเลือก และ $A$ คือเซตของ Events ที่เป็นการโจมตี

  \item \textbf{Token Cost} -- จำนวน Input Token ที่ส่งให้ LLM ในแต่ละกรณี
        และสัดส่วนจำนวน Events ที่ถูกคัดเลือกต่อ Events ทั้งหมด

  \item \textbf{Latency} -- เวลาตั้งแต่รับ Logs จนได้คำตอบ (End-to-end)
\end{itemize}

เนื่องจากจำนวนกรณีทดสอบมีจำกัด ค่าความถูกต้องจะรายงานพร้อมช่วงความเชื่อมั่น 95\%
ด้วยวิธี Bootstrap~\cite{efron1979} และเปรียบเทียบแต่ละรูปแบบการทดลองกับ Baseline
ด้วย McNemar's Test~\cite{mcnemar1947} ซึ่งเหมาะกับผลแบบถูก/ผิดที่ทดสอบบนชุดกรณีเดียวกัน
โดยปรับค่า $p$ ด้วยวิธีของ Holm~\cite{holm1979} เมื่อทดสอบหลายคู่พร้อมกัน

\section{\ifenglish%
\ifcpe CPE \else ISNE \fi knowledge used, applied, or integrated in this project
\else%
ความรู้ตามหลักสูตรซึ่งถูกนำมาใช้หรือบูรณาการในโครงงาน
\fi
}
\begin{itemize}
  \item \textbf{ความมั่นคงปลอดภัยของระบบคอมพิวเตอร์และเครือข่าย} -- ความเข้าใจพฤติกรรมการโจมตี
        Logs ด้านความปลอดภัย และกระบวนการสืบสวนเหตุการณ์
  \item \textbf{ระบบปฏิบัติการ} -- กลไกการสร้างและความสัมพันธ์ระหว่าง Process บน Windows และ Linux
  \item \textbf{โครงสร้างข้อมูลและอัลกอริทึม} -- การสร้างกราฟ การท่องกราฟด้วยการหาเส้นทางสั้นที่สุด (Dijkstra)
        และการจัดลำดับด้วยคะแนนภายใต้งบประมาณ Token
  \item \textbf{ฐานข้อมูลและการจัดการข้อมูล} -- การ Parse, Normalize และจัดเก็บข้อมูลหลายรูปแบบ
  \item \textbf{วิศวกรรมซอฟต์แวร์} -- การออกแบบสถาปัตยกรรมระบบ การแบ่งโมดูล และการพัฒนา Web Application
  \item \textbf{ปัญญาประดิษฐ์/การเรียนรู้ของเครื่อง} -- การประเมินผลด้วย Accuracy และ Recall การทดสอบนัยสำคัญทางสถิติ
        และการแยกข้อมูลสำหรับพัฒนากับข้อมูลสำหรับทดสอบ
\end{itemize}

\section{\ifenglish%
Extracurricular knowledge used, applied, or integrated in this project
\else%
ความรู้นอกหลักสูตรซึ่งถูกนำมาใช้หรือบูรณาการในโครงงาน
\fi
}
\begin{itemize}
  \item การติดตั้งและใช้งาน Splunk Enterprise รวมถึง SPL และ REST API
  \item รูปแบบ Logs ของ Windows Sysmon และ Linux Auditd และมาตรฐาน CIM
  \item กรอบความรู้ MITRE ATT\&CK และรูปแบบข้อมูล STIX
  \item การใช้งาน LLM ภายในเครื่องผ่าน Ollama, Prompt Engineering และ Structured Output
  \item การสร้างระบบ RAG ด้วย Embedding Model และ Vector Search
\end{itemize}
